\documentclass[10pt,a4paper]{amsart}
\usepackage[margin=19mm]{geometry}
\usepackage{amsmath,amssymb,mathtools}
\usepackage[scr=rsfs,cal=euler]{mathalfa}
\usepackage{xfrac}
\usepackage[unicode,hidelinks,bookmarks=false]{hyperref}
\newcommand{\ie}{i.e.\ }
\newcommand{\cf}{cf.\ }
\newcommand{\resp}{respectively\ }
\newcommand{\Subsection}{\subsection}
\providecommand{\nobreakdash}{\nobreak\hbox{-}}
\setlength{\emergencystretch}{2em}
\multlinegap0pt
\usepackage{bm}
\usepackage{bbm}
\usepackage{dsfont}
\usepackage{xspace}
\usepackage{cancel}
\usepackage{mathtools}
\usepackage{amsthm}
\makeatletter
\@ifundefined{theorem}{\newtheorem{theorem}{Theorem}}{}
\@ifundefined{proposition}{\newtheorem{proposition}[theorem]{Proposition}}{}
\makeatother
\newcommand{\hsum}[1]{\widehat{#1}}
\title[The greedy 3-sumfree sequence $S\_{1,g,g+1}$]{The greedy 3-sumfree sequence $S\_{1,g,g+1}$}
\author{Orion Shtrezi}
\date{Source: arXiv:2606.17447v1 (CC BY 4.0)}
\begin{document}
\maketitle

\noindent\emph{Source and licence.} The greedy 3-sumfree sequence $S_{1,g,g+1}$. Version arXiv:2606.17447v1 (CC BY 4.0), distributed under the Creative Commons Attribution 4.0 International licence, as recorded in the arXiv licence metadata for this exact version and shown on the abstract page https://arxiv.org/abs/2606.17447v1. Full rights notice: licenses/arxiv-2606.17447-CC-BY-4.0.txt.\par\smallskip

\noindent\textit{Editorial scope.} Counted unit: the Proposition whose source label is \texttt{prop:omitted} in the article (source0.tex lines 102--104, source label \texttt{prop:omitted}), delivered together with the article's own complete original proof of it (source0.tex lines 112--165, one \texttt{proof} environment of 54 lines). No mathematics is written, repaired or re-derived: statement and proof are the article's own text line for line. The proof is detached from the statement in the source: the article states the result at lines 102--104 and prints its proof at lines 112--165, and the proof environment's own header names the statement, so the association is the article's own. Required context retained uncounted: eq:hat-identities (lines 67--71). Every definition, display and label of those lines is retained unchanged. Omitted from this package and not redistributed: 1--66, 72--101, 105--111, 166--305. Nothing inside a retained excerpt is omitted or altered: every retained body is delivered exactly as the article's lines are written, so no omission receipt of any kind accompanies this package. Citations: the retained excerpts carry no citation command at all, so no bibliography entry of the article is transcribed and no reference is redistributed. Reference adaptation: none was needed and none was made; every reference inside the delivered unit resolves inside it, so no unresolved reference or citation is produced. Printed slips kept literal: no printed word, symbol, hypothesis or number of the counted statement or of its proof was altered. Layout and class: the article's own class is \texttt{article} with packages including amsmath, amsthm, amssymb, hyperref; the wrapper is the lane's pinned \texttt{amsart} editorial wrapper at 10pt on A4, which restates the theorem-like environments the retained text needs and adds the article's own macro definitions that the retained text needs (\texttt{\textbackslash hsum}), with extra wrapper packages (bm, bbm, dsfont, xspace, cancel, mathtools). The delivered numbering is the unit's own. Assets: no retained span contains a figure environment, a graphics-inclusion command, a PDF-page inclusion, a table or a TikZ picture; no image file is copied into this package, the package declares no asset dependency and no image file is redistributed with it. Licence: version arXiv:2606.17447v1 of this article is distributed under the Creative Commons Attribution 4.0 International licence, as recorded in the arXiv licence metadata for this exact version and shown on the abstract page https://arxiv.org/abs/2606.17447v1. A full rights notice is delivered at licenses/arxiv-2606.17447-CC-BY-4.0.txt. Characters: every retained excerpt is pure ASCII, so the delivered engine, XeLaTeX with its default Latin Modern text font, reports no missing character.
\begin{equation}\label{eq:hat-identities}
\hsum{2}[u,v]=[2u+1,2v-1],
\qquad
\hsum{3}[u,v]=[3u+3,3v-3],
\end{equation}

\begin{proposition}\label{prop:omitted}
Every integer $n>g+1$ with $n\notin A_g$ is forced.
\end{proposition}

\begin{proof}[Proof of Proposition~\ref{prop:omitted}]
For the initial block, let
\[
B=[g,2g+1],
\qquad
D=[6g+1,7g+1],
\]
so that $B=I_0\cup\{e\}$ and $D=\{f\}\cup J_0$ are subsets of $A_g$.
The omitted integers above $g+1$ in block $0$ are
\[
[2g+2,6g]\cup[7g+2,m].
\]
By \eqref{eq:hat-identities} we have
\[
1+\hsum{2}B=[2g+2,4g+2],
\qquad
\hsum{3}B=[3g+3,6g],
\]
so these two intervals cover $[2g+2,6g]$. Likewise,
\[
1+B+D=[7g+2,9g+3],
\qquad
(7g+1)+\hsum{2}B=[9g+2,11g+2].
\]
Since $m=10g+3\le 11g+2$, these two intervals cover $[7g+2,m]$.

Now let $k\ge 1$. The omitted integers in block $k$ are exactly
\[
km+[1,g-1],
\qquad
km+[2g+1,6g+1],
\qquad
km+[7g+2,m].
\]
For the first gap we have
\[
((k-1)m+7g+1)+e+[g+2,2g]=km+[1,g-1].
\]
For the middle gap we have
\[
1+B+I_k=km+[2g+1,4g+2],
\qquad
(km+2g)+\hsum{2}B=km+[4g+1,6g+1].
\]
For the last gap we have
\[
1+D+I_k=km+[7g+2,9g+2],
\qquad
(km+7g+1)+\hsum{2}B=km+[9g+2,11g+2].
\]
Again $m=10g+3\le 11g+2$, so these intervals cover all of $km+[7g+2,m]$.

Each displayed identity gives a representation by three distinct elements of $A_g$. The operators $\hsum{2}$ and $\hsum{3}$ enforce distinctness inside $B$, and the remaining singled-out summand lies outside the relevant copy of $B$. In the mixed sums $1+B+D$, $1+B+I_k$, and $1+D+I_k$, the three summands come from pairwise disjoint sets, and the same holds for the representation of the first gap, whose summands come from $J_{k-1}$, $\{e\}$, and $[g+2,2g]\subseteq I_0$. Therefore every omitted integer $n>g+1$ is forced.
\end{proof}

\end{document}
